% =====================================================================
%  ch11.tex —— 第 11 章 自由费米子可逆相与 K 理论
%  底本：原书 p.152–p.170（章标题起于 p.152 中部；p.170 顶部尚有本章
%  §11.6.2 的余下内容与式 (11.53)，其后才是第 12 章标题）
%  §11.1 两个基本例子 §11.2 十重对称类 §11.3 对称类的“移位”
%  §11.4 $0+1$ 维的有能隙哈密顿量 §11.5 $K$ 与 $KO$ 理论
%  §11.6 $d+1$ 维的有能隙哈密顿量
%  本章表 7（p.159）、表 8（p.160）、表 9（p.169）均以 tabular 重排。
% =====================================================================
\chapter{自由费米子可逆相与 $K$ 理论}

% [原文 p.152]
\section{两个基本例子}

\subsection{$1+1$ 维边界、$2+1$ 维体，无额外不变性}

\para{边界分析：}考虑 $1+1$ 维中一个无质量（massless）且无能隙（gapless）的相对论性实费米子。它既可以向左运动，
\begin{equation}
\frac{\partial}{\partial t}\psi+\frac{\partial}{\partial x}\psi=0, \tag{11.1}
\end{equation}
也可以向右运动
\begin{equation}
\frac{\partial}{\partial t}\psi-\frac{\partial}{\partial x}\psi=0. \tag{11.2}
\end{equation}
它们的色散关系为 $E=\pm p$。量子化之后，$b_p$ 与 $b_{-p}$ 成为能量为 $|p|$ 的模式的费米子产生与湮灭算子对。在凝聚态中，当某处波数上有能带穿过费米能量时就会出现这种情况；此时我们测量的是相对于该点的能量与动量，并且取时间与长度的单位使速度变为 1。在以上的一般性讨论中，它们对应于 $\mathrm{Cliff}_{\R}(1)$ 的表示，其中式 (10.40) 的 $\alpha^1$ 由 $\pm1$ 表示。

要让它们有质量（有能隙），需要把两者结合起来，得到色散关系
\begin{equation}
E=\pm\sqrt{p^2+m^2}. \tag{11.3}
\end{equation}
这来自实二分量 $\psi$ 的方程
\begin{equation}
\frac{\partial}{\partial t}\psi+\alpha_1\frac{\partial}{\partial x}\psi+m\beta\psi=0, \tag{11.4}
\end{equation}
% [原文 p.153]
其中 $\beta$ 是某个矩阵（采用非相对论性约定）。正如 Dirac 最初所论证的，为使 $ip\alpha_1+m\beta$ 具有本征值 (11.3)，我们需要
\begin{equation}
(\alpha_1)^2=+1,\qquad \beta^2=-1,\qquad \alpha_1\beta+\beta\alpha_1=0. \tag{11.5}
\end{equation}
因此，2 维中的有质量实费米子由 $\mathrm{Cliff}_{\R}(1,1)$ 的一个（不分次）表示 $\R^2$ 来刻画，正如上表所示。当我们关掉质量（能隙）时，这个表示分裂为 $\mathrm{Cliff}_{\R}(1)$ 的 $\R_+$ 与 $\R_-$。

于是，若我们有 $n_L$ 个左移模式与 $n_R$ 个右移模式，就可以把其中 $\min(n_L,n_R)$ 个配成对并使它们有质量（有能隙），而 $|n_L-n_R|$ 个仍然保持无质量（无能隙）。因此，所谓“本质上无能隙的自由度”的数目由 $\Z$ 计量。

\para{体分析：}接下来考虑 $2+1$ 维体中的一个有质量费米子，其运动方程为
\begin{equation}
\frac{\partial}{\partial t}\psi+\alpha_1\frac{\partial}{\partial x}\psi+\alpha_2\frac{\partial}{\partial y}\psi+m\beta\psi=0, \tag{11.6}
\end{equation}
其中 $\alpha_1,\alpha_2$ 与 $\beta$ 需要构成 $\mathrm{Cliff}_{\R}(2,1)$ 的一个不分次表示，它仍然是 $\R^2$。为具体起见，取
\begin{equation}
\alpha_1=\begin{pmatrix}1&0\\0&-1\end{pmatrix},\qquad
\alpha_2=\begin{pmatrix}0&1\\1&0\end{pmatrix},\qquad
\beta=\begin{pmatrix}0&-1\\1&0\end{pmatrix}. \tag{11.7}
\end{equation}
此时 $m$ 是实参数。我们有色散关系
\begin{equation}
E=\pm\sqrt{p_x^2+p_y^2+m^2}. \tag{11.8}
\end{equation}
再让质量参数 $m$ 依赖于 $y$，于是有
\begin{equation}
\frac{\partial}{\partial t}\psi+\alpha_1\frac{\partial}{\partial x}\psi+\alpha_2\frac{\partial}{\partial y}\psi+m(y)\beta\psi=0. \tag{11.9}
\end{equation}
设
\begin{equation}
\psi(t,x,y)=\Psi(t,x)f(y). \tag{11.10}
\end{equation}
那么应当有
\begin{equation}
\alpha_2\partial_y f+m(y)\beta\psi=0, \tag{11.11}
\end{equation}
或等价地
\begin{equation}
\begin{pmatrix}\partial_y+m(y)&0\\0&\partial_y-m(y)\end{pmatrix}\psi=0, \tag{11.12}
\end{equation}
也就是说
\begin{equation}
\psi=\begin{pmatrix}\exp\!\big(-\int^y m(y)dy\big)\Psi_+(t,x)\\[2pt]
\exp\!\big(+\int^y m(y)dy\big)\Psi_-(t,x)\end{pmatrix}, \tag{11.13}
\end{equation}
其中
\begin{equation}
(\partial_t+\partial_x)\Psi_+=0,\qquad (\partial_t-\partial_x)\Psi_-=0. \tag{11.14}
\end{equation}
% [Tutor Note: 式 (11.11) 原文右端作 $m(y)\beta\psi$，按 (11.10) 的分解应为 $m(y)\beta f$；此处照录原文，已登记进原文笔误表。]

% [原文 p.154]
当 $m(y)$ 在 $y\ll0$ 处取 $-M$、在 $y\gg0$ 处取 $+M$ 进行插值时，左移的无质量分量 $\Psi_+(t,x)$ 指数局域在 $y\sim0$ 附近。类似地，当 $m(y)$ 随 $y$ 从 $y\ll0$ 到 $y\gg0$ 由 $+M$ 插值到 $-M$ 时，右移的无质量分量 $\Psi_-(t,x)$ 指数局域在 $y\sim0$ 附近。

在体有能隙理论的低能极限下，无论 $m\gg0$ 还是 $m\ll0$，都不存在体激发，体几乎是完全空的。但若在 $y<0$ 处相为 $m\ll0$、在 $y>0$ 处相为 $m\gg0$，就会出现一个具有确定手征性（运动方向）的无质量费米子。我们后面会知道，体的低能极限由引力 Chern--Simons 项给出有效作用量，从而给出非零的热霍尔电导。这个非零的热霍尔电导是整数量，它计量边界上“本质上无质量自由度”的数目。

\subsection{$1+1$ 维边界、$2+1$ 维体，带 $T^2=(-1)^F$}

\para{边界分析：}让我们重复上面的分析，但这次加上时间反演不变性。在 $1+1$ 维中，时间反演把一个沿某方向运动的模式送到沿相反方向运动的模式。在无能隙（无质量）情形，最小的情形是
\begin{equation}
\gamma_0\partial_t\psi+\gamma_1\partial_x\psi=0, \tag{11.15}
\end{equation}
其中 $\psi$ 是二分量的实“pinor”；这里
\begin{equation}
\gamma_0=\begin{pmatrix}0&1\\1&0\end{pmatrix},\qquad
\gamma_1=\begin{pmatrix}0&-1\\1&0\end{pmatrix} \tag{11.16}
\end{equation}
对应 $T^2=+1$ 的 pin$^-$，或者
\begin{equation}
\gamma_0=\begin{pmatrix}0&-1\\1&0\end{pmatrix},\qquad
\gamma_1=\begin{pmatrix}0&1\\1&0\end{pmatrix} \tag{11.17}
\end{equation}
对应 $T^2=(-1)^F$ 的 pin$^+$。无论哪种情形，时间反演变换都由下式给出：
\begin{equation}
\psi'(t,x):=\gamma_0\psi(-t,x). \tag{11.18}
\end{equation}
现在假设你想让它有质量（有能隙）。我们写
\begin{equation}
\gamma_0\partial_t\psi+\gamma_1\partial_x\psi+m\gamma_m\psi=0, \tag{11.19}
\end{equation}
其中 $\gamma_m$ 是某个实矩阵。为使这个方程在式 (11.18) 下不变，$\beta$ 需要与 $\gamma_0$ 反对易。另外，与式 (11.4) 比较可知，我们需要 $\alpha:=\gamma_0\gamma_1$ 与 $\beta:=\gamma_0\gamma_m$ 反对易。这就要求 $\gamma_m$ 同时与 $\gamma_0$ 和 $\gamma_1$ 反对易。基本上唯一可能的选择就是取 $\gamma_m=\sigma_3$。% [Tutor Note: 此处原文作“$\beta$ 需要与 $\gamma_0$ 反对易”，按上下文应为 $\gamma_m$；照录原文，已登记进原文笔误表。]

此外，我们还需要 $(\gamma_0\gamma_1)^2=+1$ 与 $(\gamma_0\gamma_m)^2=-1$，这一点前面也已经论证过。后一个条件对式 (11.16) 成立，但对式 (11.17) 不成立。这意味着，对式 (11.16) 即 $T^2=+1$ 的 pin$^-$ 而言，这对左移费米子与右移费米子可以被赋予质量；而对式 (11.17) 即 $T^2=(-1)^F$ 的 pin$^+$ 而言，这是不可能的。

% [原文 p.155]
让我们集中讨论非平凡情形 $T^2=(-1)^F$。与没有时间反演的情形相比，区别在于：若有两对左移费米子与右移费米子，就可以使它们有质量。事实上，两对可以看成一个复二分量向量；于是直接取 $\gamma_m:=i\sigma_3$ 即可。因此，此情形下“本质上无质量自由度”的数目以 $\Z_2$ 而非先前的 $\Z$ 计量。

\para{体分析：}接下来考虑 $2+1$ 维中具有 $T^2=(-1)^F$ 时间反演对称性的体理论。方程为
\begin{equation}
\gamma_0\partial_t\psi+\gamma_1\partial_x\psi+\gamma_2\partial_y\psi+m\gamma_m\psi=0, \tag{11.20}
\end{equation}
其中 $\gamma_{0,1,2,m}$ 彼此反对易，并且
\begin{equation}
(\gamma_0)^2=(\gamma_m)^2=-1,\qquad (\gamma_1)^2=(\gamma_2)^2=+1. \tag{11.21}
\end{equation}
这是 $\mathrm{Cliff}_{\R}(2,2)$，其最小表示是 $\R^4$。我们想照上面那样让 $m$ 沿 $y$ 变化。当 $m$ 在 $y\sim0$ 附近由正变负时，与上面类似的计算表明，会出现两个指数局域在 $y\sim0$ 附近的无质量（无能隙）模式，一个左移、一个右移。

若取这样的系统的两份拷贝，会怎样？我们需要更细致地分析质量项的空间。先回到无质量情形。去掉 $\gamma_m$ 后，$\gamma_{0,1,2}$ 构成的是 $\mathrm{Cliff}_{\R}(1,2)$，其最小表示为 $\R^2_+$ 与 $\R^2_-$。而 $\gamma_m$ 交换这两个因子。因此我们的表示形如
\begin{equation}
\gamma_{0,1,2}=\begin{pmatrix}\square&0\\0&-\square\end{pmatrix},\qquad
m\gamma_m=\begin{pmatrix}0&-m\mathbf{1}_2\\ m\mathbf{1}_2&0\end{pmatrix}. \tag{11.22}
\end{equation}
现在我们引入这种费米子的 $n$ 份拷贝，运动方程为
\begin{equation}
\gamma_0\partial_t\psi+\gamma_1\partial_x\psi+\gamma_2\partial_y\psi+\gamma_m(\mathsf{m})\psi=0, \tag{11.23}
\end{equation}
其中 $\psi$ 有 $4n$ 个实分量，$\mathsf{m}$ 是质量参数矩阵。上面出现的矩阵可以显式取为
\begin{equation}
\gamma_{0,1,2}=\begin{pmatrix}\square\otimes\mathbf{1}_n&0\\0&-\square\otimes\mathbf{1}_n\end{pmatrix},\qquad
\gamma_m(\mathsf{m})=\begin{pmatrix}0&-\mathbf{1}_2\otimes\mathsf{m}\\ \mathbf{1}_2\otimes\mathsf{m}^T&0\end{pmatrix}, \tag{11.24}
\end{equation}
其中 $\mathsf{m}$ 现在是 $n\times n$ 实矩阵。于是色散关系为
\begin{equation}
E=\pm\sqrt{p_x^2+p_y^2+\mathsf{m}^T\mathsf{m}}. \tag{11.25}
\end{equation}
那么，当我们让 $\mathsf{m}$ 在 $y\sim0$ 附近从 $-\diag(m,\ldots,m)$ 变到 $+\diag(m,\ldots,m)$ 时，会发生什么？若选取 $y$ 依赖为
\begin{equation}
\mathsf{m}=m(y)\diag(1,1,\ldots,1) \tag{11.26}
\end{equation}
% [原文 p.156]
其中 $m(y)$ 照旧从 $-M$ 插值到 $+M$，我们就只是得到 $n$ 个指数局域在 $y\sim0$ 附近的无质量时间反演不变费米子。

然而，当 $n=2k$ 为偶数时，可以把这个路径连续形变为
\begin{equation}
\mathsf{m}(y)=M\begin{pmatrix}\cos\theta(y)&-\sin\theta(y)\\[2pt] \sin\theta(y)&\cos\theta(y)\end{pmatrix}\otimes\mathbf{1}_k, \tag{11.27}
\end{equation}
使得 $\theta(y)$ 从 $-\pi$ 变到 $+\pi$。此时体质量从不消失，因此 $y\sim0$ 附近也不出现局域模。

当 $n$ 为奇数时，不存在使体质量永不消失的这类形变。这是因为此时 $\det\mathsf{m}$ 由负变正，因而必在某处过零，从而在 $y\sim0$ 附近产生至少一对局域模。这样我们又看到，边界上“本质上无质量（无能隙）自由度”的数目由 $\Z_2$ 计量；这一事实也可以从体中 $\det\mathsf{m}$ 的符号看出来。

\takeaway{两个基本例子：无额外不变性时边界无能隙自由度按 $\Z$ 计数，体由引力 Chern--Simons 项给出整数热霍尔电导；加入 $T^2=(-1)^F$ 后改按 $\Z_2$ 计数，对应体中 $\det\mathsf{m}$ 的符号。}

\section{十重对称类}

我们希望仿照上面看到的基本例子，对自由费米子可逆相作一般性分析。在此之前，需要先研究一个乍看并不相关的课题。

过去 Wigner 曾提出，研究一个非常复杂的量子力学系统时，可以假设哈密顿量 $H$ 是随机的：其每个矩阵元在给定的系综（例如高斯系综）中随机选取；然后研究这种哈密顿量本征值的性质。一个相关的问题是研究随机酉算子 $U$ 的对角化问题。

关于这种随机矩阵，我们可以说很多，并且有一个称为随机矩阵论的大研究领域。已知它能相当好地重现复杂原子核的能级分布，这是最初预期的应用之一。不仅如此，人们还知道 Riemann zeta 函数的零点服从这类随机矩阵的同一分布，这引发了“这个数论问题也许与量子力学有关”的猜想。

在朴素的随机矩阵论中，我们考虑任意 Hermite 矩阵 $H$，或任意酉矩阵 $U$。施加对称性当然会影响它们的性质。因此，对对称性类型进行分类十分重要。下面十个对称类可以从多种不同方式加以说明。我们按历史顺序展开。

\subsection{“经典”对称类}

最基本的情形是考虑所有酉矩阵 $U$，或所有 Hermite 矩阵 $H$。注意 $U=e^{iH}$，即可以把 $H$ 看作在参数化 $U$ 的空间中单位元（原点）附近的邻域。

我们也可以要求 $U$ 或 $H$ 的矩阵元属于 $\R$ 或属于 $\mathbb{H}$。等价地说，$U$ 或 $H$ 作用其上的 希尔伯特空间 $\mathcal{H}$ 带有时间反演对称性 $T$，它是一个
% [原文 p.157]
反酉算子，满足 $T^2=+1$ 或 $T^2=-1$，并要求 $H$ 与这个 $T$ 交换：
\begin{equation}
TH=HT. \tag{11.28}
\end{equation}
由于 $U=e^{iH}$，我们有
\begin{equation}
TU=U^{-1}T. \tag{11.29}
\end{equation}
总共得到三个基本对称类。它们有相应的随机矩阵论，或等价的矩阵模型。最基本的情形是考虑这类矩阵的高斯系综，其配分函数为
\begin{equation}
\int dH\,e^{-c\tr H^\dagger H}, \tag{11.30}
\end{equation}
其中 $H$ 是 $N\times N$ 的 Hermite 矩阵，矩阵元分别在 $\C$、$\R$ 或 $\mathbb{H}$ 中，方差为 $c$。关于如何取大 $N$ 极限，有一个非常有趣的故事，但这不是我们的主题。这些大 $N$ 极限分别称为高斯酉系综、高斯正交系综与高斯辛系综（GUE、GOE、GSE）。

\subsection{“手征”对称类}

1990 年代初，Verbaarschot 等人把随机矩阵模型应用于 QCD 的研究。那里要研究的大矩阵是 Dirac 算子本身，即 $\gamma^\mu D_\mu$；它作用在夸克波函数的空间上，而该空间是无限维的，在数值格点模拟中我们用巨大但有限尺寸的矩阵来近似它。$\gamma^\mu D_\mu$ 的一个显著特征是它与 $\gamma^5$ 反对易：
\begin{equation}
\gamma^5(\gamma^\mu D_\mu)+(\gamma^\mu D_\mu)\gamma^5=0,\qquad (\gamma^5)^2=+1. \tag{11.31}
\end{equation}
把这些信息抽象出来：令 $H$ 为 Hermite 哈密顿量（上例中的 $\gamma^\mu D_\mu$），$S$ 为满足 $S^2=+1$ 的酉算子（上例中的 $\gamma^5$），二者相互反对易，$SH+HS=0$。先把 $S$ 对角化，使其具有形式
\begin{equation}
S=\begin{pmatrix}\mathbf{1}_N&0\\0&-\mathbf{1}_M\end{pmatrix}, \tag{11.32}
\end{equation}
则 $H$ 具有形式
\begin{equation}
H=\begin{pmatrix}0&h\\ h^\dagger&0\end{pmatrix}, \tag{11.33}
\end{equation}
其中 $h$ 是任意的复 $N\times M$ 矩阵。因此这个系综也可以简单地看作不带 Hermite 条件的复矩阵系综。由于 $\gamma^5$ 生成的对称性在 QCD 中称为手征对称性，这个系综称为手征 GUE。

我们可以像前面那样引入满足 $T^2=+1$ 或 $T^2=-1$、且与 $S$ 交换的反酉算子 $T$。它们等价于不带进一步条件的实矩阵或四元数 $N\times M$ 矩阵系综，分别称为手征 GOE 与手征 GSE。

% [原文 p.158]
\subsection{“BdG”或“AZ”对称类}

最后，Altland 与 Zirnbauer 在 [AZ96] 中注意到，在超导体的 Bogoliubov--de Gennes 描述中会自然出现另外四个对称类。在此语境下，考虑作用在费米子单粒子 希尔伯特空间 $\mathcal{H}$ 上的 Hermite 算子 $H$，并考察可以施加的一般对称性约束。在 [AZ96] 之前未被着重处理的是 $\mathcal{H}$ 上的反酉算子 $\tilde T$，它同样满足 $\tilde T^2=+1$ 或 $\tilde T^2=-1$，并且与 $H$ 反对易：
\begin{equation}
\tilde TH=-H\tilde T. \tag{11.34}
\end{equation}
这个关系把每个正本征值与一个负本征值配对，容易使 $H$ 无下界。但这没有关系，因为负本征态被看作空穴。所以这个 $\tilde T$ 是一类粒子--空穴对称性。

我们还可以假设存在与 $H$ 反对易的酉算子 $S$，并且 $S\tilde T=\pm\tilde T S$。当这个符号为 $+$ 时，结果只是得到手征正交或手征辛对称类，因为 $T:=\tilde TS$ 与 $H$ 和 $S$ 都交换。

因此，Altland 与 Zirnbauer 引入的“新”类就是以下四种情形：
\begin{enumerate}
  \item $\tilde T$ 满足 $\tilde T^2=+1$；
  \item $\tilde T$ 满足 $\tilde T^2=-1$；
  \item 情形 1 且 $S^2=1$、$S$ 与 $\tilde T$ 反对易；
  \item 情形 2 且 $S^2=1$、$S$ 与 $\tilde T$ 反对易。
\end{enumerate}

\subsection{小结}

总共有十个对称类。下面更详细地研究这十种情形。

\subsection{十重对称类的 Cartan 名称}

迄今引入的十个类汇总在表 7 与表 8 中。它们的命名依据 Cartan 分类，这一点最早由 Altland 与 Zirnbauer 在上述论文中指出。解释如下。

在每个对称类中，我们用 $S$、$T$ 与 $\tilde T$ 的某个子集来限制酉算子 $U$，条件如下：
\begin{itemize}
  \item $S$ 是酉的，$S^2=1$，且 $SU=U^{-1}S$，对应 $SH=-HS$。
  \item $T$ 是反酉的，$T^2=\pm1$，且 $TU=U^{-1}T$，对应 $TH=+HT$。
  \item $\tilde T$ 是反酉的，$\tilde T^2=\pm1$，且 $\tilde TU=U\tilde T$，对应 $\tilde H=-H\tilde T$。
\end{itemize}
这里在 $U=e^{itH}$ 时也列出了对 $H$ 的条件。% [Tutor Note: 第三条原文作 $\tilde H=-H\tilde T$，按上下文应为 $\tilde TH=-H\tilde T$；照录原文，已登记进原文笔误表。]

下面考虑在选定的一组约束下所有这种酉算子 $U$ 构成的空间。我们逐一来看。设 $U$ 作用在 $\C^N$ 上。

% [原文 p.159]
\begin{center}
\small
\begin{tabular}{c||c|c|c|c|c|c}
\toprule
$i$ & grading & name & unitaries & \begin{tabular}{c}gapped\\[-2pt]Hamiltonians\end{tabular} & $K^{-i}(\pt)$ & hist.\\
\midrule
0 & & A & $U$ & $BU\times\Z$ & $\Z$ & classic\\
1 & \checkmark & AIII & $BU\times\Z$ & $U$ & 0 & chiral\\
\bottomrule
\end{tabular}

\smallskip
{\small 表 7：对称类：两个复类。}
\end{center}

\begin{itemize}
  \item 勾选“grading”表示存在一个固定的酉算子 $S$，满足 $S^2=1$，使得 $SU=U^{-1}S$ 或 $SH=HS$。
  \item “unitaries”与“gapped mass terms”是具有指定对称性质的酉算子 $U$ 或有能隙哈密顿量 $H$ 的分类空间。
  \item “name”是酉算子分类空间的 Cartan 名称，也用来命名对称类。
  \item $K^{-i}(\pt)$ 是有能隙哈密顿量分类空间的 $\pi_0$。
\end{itemize}

\begin{itemize}
  \item 没有任何约束时，所有这样的 $U$ 显然构成空间 $U(N)$。
  \item 当 $\tilde T^2=+1$ 时，$U$ 作用在 $\tilde T=+1$ 的实子空间 $\R^N$ 内。此时 $U$ 是实矩阵，因此它们构成空间 $O(N)$。
  \item 当 $\tilde T^2=-1$ 时，可以把 $\tilde T$ 看作四元数 $j$ 的乘法，此时 $N=2n$。于是 $U$ 是四元数酉的，构成 $Sp(n)$。
\end{itemize}
到此为止看到的都是群流形。现在把 $S$ 也包含进来。
\begin{itemize}
  \item 当 $S^2=+1$ 时，$\C^N$ 可按 $S=\pm1$ 的本征值分解为两个子空间 $V_+$ 与 $V_-$。现在考虑 $(US)^2=USUS=+1$。因此 $US$ 可用来把 $\C^N$ 分成 $W_+$ 与 $W_-$，其中 $W_\pm$ 是 $US$ 的本征值为 $\pm1$ 的本征空间。反过来，$W_+$ 与 $W_-$ 这一对决定了 $US$，从而决定了 $U$。所以 $U$ 的空间就是子空间对 $(W_+,W_-)$ 的空间。更具体些，设 $u$ 是把 $V_+$ 映到 $W_+$、把 $V_-$ 映到 $W_-$ 的酉矩阵。则 $US=uSu^{-1}$，且 $U$ 由 $u$ 决定。这个 $u$ 并不唯一：当 $u_0$ 与 $S$ 交换时，$u$ 与 $uu_0$ 给出同一个 $U$。这样的 $u_0$ 构成群 $U(N_+)\times U(N_-)$，其中 $N_\pm$ 是 $V_\pm$ 的维数。因此 $U$ 的空间就是商空间 $U(N)/(U(N_+)\times U(N_-))$。
  \item 当 $S^2=+1$ 且 $\tilde T^2=+1$、$\tilde T$ 与 $S$ 交换时，可以在 $\R^N$ 而非 $\C^N$ 上工作，得到 $U$ 的空间是商空间 $O(N)/(O(N_+)\times O(N_-))$。
  \item 当 $S^2=+1$ 且 $\tilde T^2=-1$、$\tilde T$ 与 $S$ 交换时，可以在 $\mathbb{H}^n$（$2n=N$）而非 $\C^N$ 上工作，得到 $U$ 的空间是商空间 $Sp(n)/(Sp(n_+)\times Sp(n_-))$。
\end{itemize}

% [原文 p.160]
\begin{center}
\footnotesize
\setlength{\tabcolsep}{4pt}
\begin{tabular}{c||c|c|c|c|c|c|c|c}
\toprule
$i$ & $T^2$ & $\tilde T^2$ & grading & name & unitaries & \begin{tabular}{c}gapped\\[-2pt]Hamiltonians\end{tabular} & $KO^{-i}(\pt)$ & hist.\\
\midrule
0 & $+1$ & & & AI & $U/O$ & $BO\times\Z$ & $\Z$ & classic\\
1 & $+1$ & $+1$ & $+$ & BDI & $BO\times\Z$ & $O$ & $\Z_2$ & chiral\\
2 & & $+1$ & & D & $O$ & $O/U$ & $\Z_2$ & AZ\\
3 & $-1$ & $+1$ & $-$ & DIII & $O/U$ & $U/Sp$ & 0 & AZ\\
4 & $-1$ & & & AII & $U/Sp$ & $BSp\times\Z$ & $\Z$ & classic\\
5 & $-1$ & $-1$ & $+$ & CII & $BSp\times\Z$ & $Sp$ & 0 & chiral\\
6 & & $-1$ & & C & $Sp$ & $Sp/U$ & 0 & AZ\\
7 & $+1$ & $-1$ & $-$ & CI & $Sp/U$ & $U/O$ & 0 & AZ\\
\bottomrule
\end{tabular}

\smallskip
{\small 表 8：对称类：八个实类。}
\end{center}

\begin{itemize}
  \item $T$ 是反酉算子，满足 $T^2=\pm1$，且 $TU=U^{-1}T$ 或 $TH=HT$。
  \item $\tilde T$ 是反酉算子，满足 $\tilde T^2=\pm1$，且 $\tilde TU=U\tilde T$ 或 $\tilde TH=-H\tilde T$。
  \item “grading”表示存在酉算子 $S$，$S^2=1$，使得 $\tilde T=TS$；其符号表示 $TS=\pm ST$。它自动满足 $SU=U^{-1}S$ 或 $SH=HS$，与复对称类中一样。
  \item “unitaries”“gapped mass terms”是满足指定对称性质的酉算子 $U$ 或有能隙哈密顿量 $H$ 的分类空间。
  \item “name”是酉算子分类空间的 Cartan 型，也用来命名对称类。
  \item $KO^{-i}(\pt)$ 是有能隙哈密顿量分类空间的 $\pi_0$。
\end{itemize}

% [原文 p.161]
这三个空间称为 Grassmann 流形。在取极限 $N=N_++N_-\to\infty$ 并保持 $N_-$ 固定时，它们分别成为分类空间 $BU(N_-)$、$BO(N_-)$ 与 $BSp(n_-)$。

现在还剩四个涉及 $T$ 的情形尚未处理。
\begin{itemize}
  \item 当 $T^2=+1$ 时，有 $TU=U^{-1}T$。回忆 $U^{-1}=U^\dagger$，并在 $T$ 就是矩阵元复共轭的基下工作，可知 $U=U^T$，即 $U$ 是对称酉矩阵。暂时把这种矩阵的集合记作 $\mathcal{U}$。它被 $u\in U(N)$ 的作用 $U\mapsto uUu^T$ 所保持。给定点 $\mathbf{1}_N\in U(N)$，固定它的 $U(N)$ 的子群是 $\{u\mid uu^T=\mathbf{1}_N\}$，即正交矩阵群 $O(N)$。因此 $\mathcal{U}=U(N)/O(N)$。
  \item 当 $T^2=-1$ 时，类似地可以在四元数基下工作，证明 $U$ 的空间是 $U(N)/Sp(n)$，其中 $N=2n$。
  \item 当 $\tilde T^2=+1$ 且 $S\tilde T=-\tilde TS$（因而 $T^2=-1$）时，先用 $\tilde T$ 在 $\R^N$ 中工作，于是 $U$ 被限制在 $O(N)$ 内。现在 $I:=iS$ 把 $\R^N$ 映到自身，且平方为 $-1$，从而赋予 $\R^N$ 复向量空间 $\C^n$ 的结构，其中 $N=2n$。同样有 $(UI)^2=-1$。与前面一样，$UI=uIu^{-1}$，且 $u$ 决定 $U$。此外，当 $u_0I(u_0)^{-1}=I$ 时，$u$ 与 $uu_0$ 给出同一个 $U$。这样的 $u_0$ 构成群 $U(n)$。因此 $U$ 的空间是商空间 $O(2n)/U(n)$。
  \item 最后一种情形是 $\tilde T^2=-1$、$S\tilde T=-\tilde TS$（因而 $T^2=+1$）。与上面类似地分析可得，$U$ 的空间是 $Sp(n)/U(n)$，其中 $N=2n$。
\end{itemize}
总之，我们得到以下十个酉算子空间：
\begin{equation}
\begin{gathered}
U(N),\qquad O(N),\qquad Sp(n),\\[2pt]
\frac{U(N)}{U(N_+)\times U(N_-)},\qquad
\frac{O(N)}{O(N_+)\times O(N_-)},\qquad
\frac{Sp(n)}{Sp(n_+)\times Sp(n_-)},\\[2pt]
\frac{U(N)}{O(N)},\qquad
\frac{U(N)}{Sp(n)},\qquad
\frac{O(2n)}{U(n)},\qquad
\frac{Sp(n)}{U(n)}.
\end{gathered} \tag{11.35}
\end{equation}
它们都是齐性空间（即形如 $G/H$ 的空间，$G$ 是某个群，$H$ 是子群）。不仅如此，它们还都是对称空间（即群 $G$ 在某个 2 阶作用 $\sigma:G\to G$ 下的不变子集，同时也是齐性空间 $G/H$）。Cartan 分类了这类对称空间，发现十个无限族以及若干例外情形。我们上面找到的十个族正好就是 Cartan 找到的无限族，而且穷尽了它们。Cartan 用 $G$ 的 Dynkin 名（即 $U$ 为 A、奇 $SO$ 为 B、$Sp$ 为 C、偶 $SO$ 为 D 等等）再加上罗马数字来区分分母 $H$，这些名称见表 7 与表 8。

在表中我们取了极限 $N\to\infty$，即扩大 希尔伯特空间 $\mathcal{H}$。习惯上，把极限的无限维对称空间记作去掉 $N$ 或 $n$ 的记号。对于 $G=U,O,Sp$ 的 $G/(G\times G)$，我们改用 $BG$，因为它们是 $G$-丛的分类空间。
% [原文 p.162]
此外，我们写 $BG\times\Z$ 而不只是 $BG$，是为了提醒自己 $N_+$ 与 $N_-$ 不一定相等，并且多了一个整数自由度的参数。（这个取无限维极限的过程可以更完整、更严格地描述，但这里不展开。）

\takeaway{十重对称类源于随机矩阵论：三个经典类、手征类与 BdG 类；其酉算子空间是 Cartan 的十族对称空间，有能隙哈密顿量的空间则比酉算子空间移位一步（表 7、表 8）。}

\section{对称类的“移位”}

十个类可以自然地分成两个复类与八个实类，取决于约束是否涉及反酉算子。由于复类只有两个，给它们排序似乎没有意义；但在表中我们确实为八个实类选取了一个特殊的顺序。这与费米子的性质有关，下面就来解释。

\subsection{复对称类的移位}

从复类开始。类 A 是所有酉算子 $U$ 的空间，对 $U$ 没有任何约束。

让我们加上一个约束：存在费米子算子 $\psi$，满足 $\psi U=U^{-1}\psi$。（若 $U=e^{iH}$，这等价于 $H\psi+\psi H=0$。）这样的 $U$ 的空间是什么？其实 $\psi$ 只是 $S$ 的另一个名字，于是得到对称类 AIII。

如果再加入两个相互反对易、并满足 $\psi_iU=U^{-1}\psi_i$ 的费米子算子 $\psi_1$ 与 $\psi_2$，会怎样？换言之，若给对称类 AIII 再加上一个满足 $S'^2=1$、$S'S=-SS'$ 的酉算子 $S'$，会怎样？设 $\mathcal{H}$ 是 $U$ 作用其上的 希尔伯特空间。算子 $iSS'$ 的平方为 1。令 $\mathcal{H}_+$ 为 $iSS'=+1$ 的本征空间。由于 $iSS'$ 与 $U$ 交换，我们只是在 $\mathcal{H}_+$ 上有一个酉算子 $U$，于是回到了对称类 A。

\subsection{实对称类的移位}

现在考虑实类。从类 AI 出发，$TU=U^{-1}T$ 且 $T^2=+1$。若 $U=e^{iH}$，这等价于 $TH+HT=0$。

如果加入相互反对易、并满足 $\psi_iU=U^{-1}\psi_i$ 的费米子算子 $\psi_{1,2,\ldots,n}$，会怎样？此时还需要指定 $\psi_i$ 如何与 $T$ 交换。我们考虑 $\psi_iT=T\psi_i$ 的情形。
\begin{enumerate}
  \item 只有一个 $\psi$，即 $n=1$：直接把 $\psi$ 改名为 $S$，就得到类 BDI。
  \item 有两个费米子算子 $\psi_1$ 与 $\psi_2$：考虑平方为 1 的算子 $i\psi_1\psi_2$。令 $\mathcal{H}_+$ 为 $i\psi_1\psi_2=+1$ 的本征空间。$U$ 在 $\mathcal{H}_+$ 内作用，但 $T$ 不作用。改为考虑 $T':=T\psi_1$。它保持 $\mathcal{H}_+$，平方为 $+1$，并满足 $T'U=UT'$。因此得到类 D，其中 $\tilde T:=T'$。
  \item 有三个费米子算子 $\psi_1,\psi_2,\psi_3$：先用前两个进入类 D。此时 $\psi_3$ 与 $\tilde T$ 反对易。记 $S:=\psi_3$，这就是类 DIII。
% [原文 p.163]
  \item 有第四个费米子算子 $\psi_4$：考虑 $i\psi_3\psi_4$ 的 $+1$ 本征空间。它被 $U$ 与 $T'':=\tilde T\psi_3$ 保持。$T''$ 平方为 $-1$，并满足 $T''U=U^{-1}T''$。这就是类 AII。
  \item 有第五个费米子算子 $\psi_5$：$T''$ 与 $\psi_5$ 反对易。记 $S:=\psi_5$，就得到类 CII。
  \item 有第六个费米子算子 $\psi_6$：考虑 $i\psi_5\psi_6$ 的 $+1$ 本征空间。它被 $U$ 与 $T''':=T''\psi_5$ 保持。$T'''$ 平方为 $-1$，并满足 $T'''U=UT'''$。令 $\tilde T:=T'''$，就得到类 C。
  \item 有第七个费米子算子 $\psi_7$：$T'''$ 与 $\psi_7$ 交换。记 $S:=\psi_7$，就得到类 CI。
  \item 最后，有第八个费米子算子 $\psi_8$：考虑 $i\psi_7\psi_8$ 的 $+1$ 本征空间。它被 $U$ 与 $T'''':=T'''\psi_7$ 保持。$T''''$ 平方为 $+1$，并满足 $T''''U=U^{-1}T''''$。这就是类 AI，即回到了出发点。
\end{enumerate}
这一对称性移位过程可以无限继续，周期为模 8。在表 8 中，我们按此顺序排列各类，并从 $i=0$ 到 $i=7$ 编号。% [Tutor Note: 正文第 4 条原文的 $T''$ 与第 6 条的 $T'''$ 按移位链依次定义，注意与表 8 的 $i$ 编号对应。]

继续之前再作一点说明。设有作用在 $\mathcal{H}_1$ 上的第 $i$ 个对称类的系统，以及作用在 $\mathcal{H}_2$ 上的第 $j$ 个对称类的系统。$\mathcal{H}_1\otimes\mathcal{H}_2$ 上的组合系统从两个系统继承时间反演对称性，但二者以非平凡的方式组合。结果表明，组合系统属于第 $(i+j)$ 个对称类。可以逐一验证，但更简单的办法是注意：第 $i$ 个对称类就是 $TU=U^{-1}T$、$T^2=+1$，再加上 $i$ 个费米子算子。于是组合系统就是 $TU=U^{-1}T$、$T^2=+1$，再加上 $(i+j)$ 个费米子算子。这既显示了“移位”过程的用处，也显示了对称类相应编号的用处。

\takeaway{每加入一个费米子算子，对称类就移位一步：复类在 A 与 AIII 之间 2-周期，实类从 AI 经 BDI 等回到 AI 为 8-周期；两个系统组合后的类按 $i+j$ 相加。}

\section{$0+1$ 维的有能隙哈密顿量}

到此为止我们考虑的都是酉算子 $U$ 的对称类。现在选定一个对称类，假设 $U=e^{iH}$，并进一步要求 $H$ 没有零本征值。这样的 $H$ 称为有能隙哈密顿量。

这里的 $H$ 应当看作费米子系统中的单粒子哈密顿量。用凝聚态的语言说，能量为零对应于费米能量，正本征值是粒子，负本征值是空穴。用 hep-th 的语言说，费米能量隐含地取为零，正本征值是粒子，负本征值是反粒子。无论哪种说法，只要 $H$ 没有零本征值，系统就是有能隙（有质量）的。

我们希望把这些有能隙哈密顿量在连续形变下分类。先从 $0+1$ 维情形开始。考虑复对称类 A。有能隙哈密顿量 $H$ 的空间就是所有没有零本征值的 Hermite 矩阵 $H$ 构成的空间。考虑在保持本征空间不变的前提下改变本征值。由于不越过本征值零，我们可以把所有本征值都取为 $\pm1$。这样就单粒子 希尔伯特空间
% [原文 p.164]
$\mathcal{H}$ 分成了正、负本征空间 $\mathcal{H}_+$ 与 $\mathcal{H}_-$。这种分裂前面讨论另一个复对称类，即类 AIII 时出现过，它由 $U(N)/(U(N_+)\times U(N_-))$ 参数化。
% [Tutor Note: 原文此处作“the class AII”，按上下文应为 AIII（另一个复对称类）；已登记进原文笔误表。]对固定的分裂 $N=N_++N_-$，该空间是连通的；而两个不同的分裂 $(N_+,N_-)\neq(N'_+,N'_-)$ 彼此不连通。因此复对称类 A 中有能隙哈密顿量的空间由一个整数 $N_+\in\Z$ 分类。可以取大 $N$ 极限，此时有能隙哈密顿量的空间成为
\begin{equation}
\prod_{N_+}\frac{U(N)}{U(N_+)\times U(N_-)}\xrightarrow{N\to\infty}BU\times\Z. \tag{11.36}
\end{equation}
于是分类由下式给出：
\begin{equation}
\pi_0(BU\times\Z)=\Z. \tag{11.37}
\end{equation}
实对称类 AI 类似。唯一的区别是强加了与 $H$ 交换的时间反演对称性 $T^2=1$。此时 $\mathcal{H}$ 实际上是实向量空间，于是得到 $BO\times\Z$ 而非 $BU\times\Z$。有能隙哈密顿量的空间由下式分类：
\begin{equation}
\pi_0(BO\times\Z)=\Z. \tag{11.38}
\end{equation}
要分析其他对称类，最好系统地来做。我们介绍两种方法。

先作两种方法共同的准备。为记号简洁，假定我们研究的是实对称类；复对称类可以类似处理，而且更容易。我们可以保持本征空间不变，把 $H$ 的本征值连续形变到 $\pm1$，而不让任何本征值变成零。这不改变哈密顿量的拓扑性质\footnote[17]{更确切地说，有能隙哈密顿量的空间与满足 $H^2=1$ 的有能隙哈密顿量的空间，在这个操作下互为形变收缩。}。于是 $H^2=+1$。

\subsection{直接分析}

先考虑不含 $S$ 的对称类，即类 AI、D、AII 与 C。此时取一个第 $i$ 型的参考有能隙哈密顿量 $H_0$。那么所有第 $i$ 型的有能隙哈密顿量 $H$ 都可以写成 $H=UH_0$，其中 $U=HH_0$。于是分类所有可能的 $H$ 等价于分类所有可能的 $U$。现在，简单计算表明 $H_0U=U^{-1}H_0$。因此，令 $S:=H_0$，本质上就是在多一个 $S$ 的条件下分类 $U$。这对应第 $i\pm1$ 个对称型，取决于 $S$ 如何与第 $i$ 型中出现的 $T$ 或 $\tilde T$ 交换。逐一分析可知，所得类型总是 $i+1$。这就确定了这四种类型中有能隙哈密顿量的空间。

接下来考虑含 $S$ 的对称类，即 BDI、DIII、CII 与 CI。此时把 希尔伯特空间 $\mathcal{H}$ 分成两个子空间 $\mathcal{H}_+$ 与 $\mathcal{H}_-$，其中 $\mathcal{H}_+$ 是 $S$ 的本征值 $+1$ 的本征空间，$\mathcal{H}_-$ 是 $S$ 的本征值 $-1$ 的本征空间。由于对易关系 $SH=-HS$，$H$ 交换 $\mathcal{H}_+$ 与 $\mathcal{H}_-$。

% [原文 p.165]
取一个第 $i$ 型的参考有能隙哈密顿量 $H_0$。那么所有第 $i$ 型的有能隙哈密顿量 $H$ 都可以像前面那样写成 $H=UH_0$，其中 $U=HH_0$。这个 $U$ 分别在 $\mathcal{H}_+$ 与 $\mathcal{H}_-$ 内作用。当 $S$ 与 $T$ 交换、且 $\tilde T=TS$（即 BDI 与 CII）时，$T$ 与 $\tilde T$ 在 $\mathcal{H}_+$ 内作为同一个算子作用，记作 $T'$。按构造 $T'U=UT'$。因此它算作 $\tilde T$ 型的反酉算子，于是 BDI 与 CII 中有能隙 $H$ 的分类分别归约为 D 与 C 中 $U$ 的分类。当 $S$ 与 $T$ 反对易、且 $\tilde T=TS$（即 DIII 与 CI）时，我们改为考虑在 $\mathcal{H}_+$ 内作用的反酉算子 $T'=TH_0$；可以验证它满足 $T'U=U^{-1}T'$。因此它算作对 $T$ 型 $U$ 的约束，于是分别得到 AII 型与 AI 型 $U$ 的分类问题。

总之，在全部八种情形中，我们总是得到第 $i+1$ 型 $U$ 的分类。

\subsection{用 Clifford 代数}

还有另一种方法，基于先前对 Clifford 表示的分析。与其研究第 $n$ 型的有能隙哈密顿量 $H$，不如研究带 $\psi_{1,\ldots,n}$ 的第 $0$ 型有能隙哈密顿量 $H$。现在保持本征空间不变，把 $H$ 的本征值形变到 $\pm1$。于是 $H$ 满足 $H^2=+1$。按构造，$H$ 与 $\psi_n$ 反对易。记 $\psi_{n+1}:=H$，可见我们试图通过适当选取 $\psi_{n+1}$，把 $\psi_{1,\ldots,n}$ 的一个固定不分次 Clifford 表示延拓为 $\psi_{1,\ldots,n+1}$ 的不分次 Clifford 表示。

现在，给定一组 $\psi_{1,\ldots,n+1}$，令 $G_n$ 为使下式成立的 $U$ 构成的群：
\begin{equation}
U\psi_iU^{-1}=\psi_i,\qquad i=1,\ldots,n. \tag{11.39}
\end{equation}
那么 $\tilde\psi_{n+1}:=U\psi_{n+1}U^{-1}$ 也满足要求。但若 $V\psi_{n+1}V^{-1}=\psi_{n+1}$，即 $V$ 属于使下式成立的 $V$ 构成的群 $G_{n+1}$，则 $UV\in G_n$ 与 $U\in G_n$ 给出同一个 $\tilde\psi_{n+1}$：
\begin{equation}
V\psi_iV^{-1}=\psi_i,\qquad i=1,\ldots,n+1. \tag{11.40}
\end{equation}
于是，固定 $\psi_{1,\ldots,n}$ 时所有 $\psi_{n+1}$ 构成的空间就是商空间 $G_n/G_{n+1}$。下面把它们算出来。

为便于跟上下面的论证，请参考表 5 回忆不分次 Clifford 模的结构：
\begin{itemize}
  \item $n=0$ 时。$\mathrm{Cliff}_{\R}(1,0)$ 的不分次表示是 $N_+$ 份 $\R_+$ 与 $N_-$ 份 $\R_-$ 的直和，而 $\mathrm{Cliff}_{\R}(0,0)$ 表示在 $\R^N$ 上，其中 $N=N_++N_-$。所以 $G_0=O(N)$、$G_1=O(N_+)\times O(N_-)$，因此 $\psi_1$ 的空间是 $O(N)/(O(N_+)\times O(N_-))$。
  \item $n=1$ 时。$\mathrm{Cliff}_{\R}(2,0)$ 表示在 $N$ 份 $\R^2$ 上，分解为表示在 $N$ 份 $\R_+$ 与 $N$ 份 $\R_-$ 上的 $\mathrm{Cliff}_{\R}(1,0)$。于是 $G_1=O(N)\times O(N)$，而 $G_2=O(N)$。所以 $\psi_2$ 的空间是 $(O(N)\times O(N))/O(N)=O(N)$。
  \item $n=2$ 时。$\mathrm{Cliff}_{\R}(3,0)$ 表示在 $N$ 份 $\C^2$ 上，而在 $\mathrm{Cliff}_{\R}(2,0)$ 下它们成为 $2N$ 份 $\R^2$。因此 $G_2=O(2N)$，而 $G_3=U(N)$，$\psi_3$ 的空间是 $O(2N)/U(N)$。
  \item $n=3$ 时。$G_3=U(2N)$、$G_4=Sp(N)$，所以 $\psi_4$ 的空间是 $Sp(N)/U(2N)$。
  \item $n=4$ 时。$G_4=Sp(N)$、$G_5=Sp(N_+)\times Sp(N_-)$，所以 $\psi_5$ 的空间是 $Sp(N)/(Sp(N_+)\times Sp(N_-))$。
  \item $n=5$ 时。$G_5=Sp(N)\times Sp(N)$、$G_6=Sp(N)$，所以 $\psi_6$ 的空间是 $(Sp(N)\times Sp(N))/Sp(N)=Sp(N)$。
  \item $n=6$ 时。$G_6=Sp(N)$、$G_7=U(N)$，$\psi_7$ 的空间是 $Sp(N)/U(N)$。
  \item $n=7$ 时。$G_7=U(N)$、$G_8=O(N)$，$\psi_8$ 的空间是 $U(N)/O(N)$。
\end{itemize}

% [原文 p.166]
这样得到的各对称类有能隙哈密顿量的空间也列在表 7 与表 8 中。注意第 $i$ 型有能隙哈密顿量的空间与第 $i-1$ 型酉算子的空间相同\footnote[18]{这也意味着，人们本可以用有能隙哈密顿量的空间而非酉算子的空间来命名对称类；那样名称会整体移动一步。}。于是有能隙哈密顿量的分类由取其空间的 $\pi_0$ 给出，它们列在表 7 与表 8 的 $K^{-i}(\pt)$ 与 $KO^{-i}(\pt)$ 两列中。下面需要解释这个奇怪的记号。

\takeaway{把 $U=e^{iH}$ 中的 $H$ 形变到本征值 $\pm1$ 之后，有能隙哈密顿量的分类化为分类空间的问题：复类给出 $BU\times\Z$，实类给出 $BO\times\Z$ 等；直接分析与 Clifford 代数两种办法得到同样的“移位一步”结论。}

\section{$K$ 与 $KO$ 理论}

定义 $K$ 理论与 $KO$ 理论有许多途径，最终都可以证明彼此相同。与我们的考虑最接近的一种定义如下：
\begin{equation}
K^0(X)=[X,BU\times\Z],\qquad K^{-1}(X)=[X,U], \tag{11.41}
\end{equation}
而一般的 $K^i(X)$ 通过考虑 $i\bmod 2$ 来定义。这里 $X$ 是一个空间，$[X,Y]$ 是从 $X$ 到 $Y$ 的映射的同伦类。类似地，
\begin{equation}
\begin{gathered}
KO^0(X)=[X,BO\times\Z],\qquad KO^{-1}(X)=[X,O],\\
KO^{-2}(X)=[X,O/U],\qquad KO^{-3}(X)=[X,U/Sp],\\
KO^{-4}(X)=[X,BSp\times\Z],\qquad KO^{-5}(X)=[X,Sp],\\
KO^{-6}(X)=[X,Sp/U],\qquad KO^{-7}(X)=[X,U/O],
\end{gathered} \tag{11.42}
\end{equation}
而一般的 $KO^i(X)$ 通过考虑 $i\bmod 8$ 来定义。

注意，上面 $K^{-i}$ 与 $KO^{-i}$ 的定义中 $[X,-]$ 第二个位置出现的空间，正好就是参数化第 $i$ 个复/实对称型的有能隙哈密顿量的对称空间。当代数拓扑中某个依赖于 $X$ 的概念 $C$ 可以用映射的同伦类 $[X,C]$ 表示时，空间 $C$ 通常称为 $C$ 的分类空间。在这个意义上，上面给出的无限维对称空间就是 $K$ 理论与 $KO$ 理论的分类空间。

取 $X=\pt$，则复/实对称型 $i$ 的 $0+1$ 维有能隙哈密顿量由 $K^{-i}(\pt)$ 与 $KO^{-i}(\pt)$ 分类，正如我们上面已经看到的那样。

% [原文 p.167]
通常，$K^0(X)$ 与 $KO^0(X)$ 是用另一种定义引入的。考虑复向量丛 $E\to X$。给定两个这样的丛 $E_1\to X$ 与 $E_2\to X$，可以构造直和丛 $E_1\oplus E_2\to X$。这使得 $X$ 上复丛的等价类 $[E]$ 的集合成为一个幺半群（即带有满足结合律的加法的集合）。我们可以为 $[E]$ 引入形式加法逆元 $-[E]$，使得 $[E]+(-[E])=0$。

为具体做到这一点，回忆一下可以由 $\N$ 构造 $\Z$：考虑数对 $(a,b)\in\N^2$，加法为 $(a,b)+(c,d)=(a+c,b+d)$，等价关系为 $(a+c,b+c)\sim(a,b)$。于是 $(0,a)$ 可以看作 $-a$。类似地，我们只需考虑数对 $([E],[F])$，加法为 $([E],[F])+([E'],[F'])=([E]+[E'],[F]+[F'])$，等价关系为 $([E]+[E'],[F]+[E'])\sim([E],[F])$。于是 $(0,[E])$ 就是 $-[E]$。如果愿意，可以把数对 $([E],[F])$ 记作 $[E\ominus F]$，并说 $E\ominus F$ 是向量丛的形式差。

如此便构造了 $X$ 上复丛等价类的形式差构成的 Abel 群。这就是 $K^0(X)$ 的标准定义（当 $X$ 紧时）。把复丛换成实丛，就得到 $KO^0(X)$。我们也可以用四元数丛定义 $KSp^0(X)$，但结果它等于 $KO^4(X)$。

我们以两种方式定义了 $K^0(X)$ 与 $KO^0(X)$：一种用丛，另一种用分类空间。要证明这两种定义等价需要相当多的工作，尽管对 $\C$ 与 $\R$ 的论证完全平行。粗略的想法如下。

首先可以证明，任何 $n$ 维复向量丛 $E\to X$ 都可以实现为平凡 $N$ 维向量丛 $\C^N\times X$ 的子丛，只要 $N\gg n$。在每一点 $x\in X$ 处，$E_x$ 定义了 $\C^N$ 的一个 $n$ 维子空间，即 Grassmann 流形 $U(N)/(U(n)\times U(N-n))$ 中的一个点。保持 $n$ 固定取极限 $N\to\infty$，得到
\begin{equation}
BU(n):=\lim_{N\to\infty}\frac{U(N)}{U(n)\times U(N-n)}, \tag{11.43}
\end{equation}
而 $X$ 上 $n$ 维复向量丛的等价类与 $[X,BU(n)]$ 一一对应。

现在需要考虑形式差 $E\ominus F$。先来说明只需考虑形如 $[E]-[\C^N]$ 的形式差，其中 $[\C^N]$ 是 $X$ 上的平凡丛 $\C^N\times X$。为此利用上面已经用过的事实：任何 $F$ 都是某个平凡丛 $\C^N$ 的子丛。于是存在 $F'$，使得 $[F]+[F']=N$，这里把 $[\C^N]$ 简记为 $N$。那么 $[E]-[F]=[E\oplus F']-N$。

我们还可以把 $[E]-N$ 改写为
\begin{equation}
\big([E]-\dim E\big)+k \tag{11.44}
\end{equation}
其中 $k$ 是整数。于是任何形式差都可以用一个整数 $k\in\Z$ 与一个向量丛 $[E]$ 来描述。$n$ 维向量丛由 $BU(n)$ 分类，取 $n\to\infty$，可见任意向量丛 $E$ 由 $BU$ 分类。连同 $k$ 一起，得到 $K^0(X)\simeq[X,BU\times\Z]$。类似地可得 $KO^0(X)\simeq[X,BO\times\Z]$ 与 $KSp^0(X)\simeq[X,BSp\times\Z]$。

$K^{-i}(X)$ 与 $KO^{-i}(X)$ 还有几种看起来不同的定义。一种是用附加 Clifford 代数作用的复/实丛。另一种定义是
% [原文 p.168]
把它们定义为
\begin{equation}
K^{-i}(X):=[X,\Omega^i(BU\times\Z)],\qquad
KO^{-i}(X):=[X,\Omega^i(BO\times\Z)] \tag{11.45}
\end{equation}
其中 $i$ 为正，$\Omega Y$ 是 $Y$ 中的环路空间，且 $\Omega^iY:=\Omega^{i-1}(\Omega Y)$。可以证明
\begin{equation}
\Omega(BU\times\Z)\simeq U,\qquad \Omega U\simeq BU, \tag{11.46}
\end{equation}
类似地
\begin{equation}
\begin{gathered}
\Omega(BO\times\Z)\simeq O,\qquad \Omega O=O/U,\\
\Omega O/U=U/Sp,\qquad \Omega U/Sp=BSp\times\Z,\\
\Omega(BSp\times\Z)\simeq Sp,\qquad \Omega Sp=Sp/U,\\
\Omega Sp/U=U/O,\qquad \Omega U/O=BO\times\Z.
\end{gathered} \tag{11.47}
\end{equation}
于是上面定义的 $K^{-i}(X)$ 与 $KO^{-i}(X)$（此时只对非负 $i$ 定义）分别是 2-周期与 8-周期的，然后可以利用这一性质把 $i$ 按周期性延拓到负值。

要看出关系 (11.46) 与式 (11.47) 为何可信，一种办法是回忆：例如实对称型 $i=2$ 时，酉算子的空间是 $O$，而有能隙哈密顿量的空间是 $O/U$。现在取一个有能隙哈密顿量 $H$。由它可以通过 $e^{itH}$ 构造酉算子空间中的一条“环路”，其中 $t$ 从 $0$ 变到 $2\pi$——只要记得 $H$ 的所有本征值都已改成了 $\pm1$。这给出映射 $O/U\to\Omega O$。其他对称型也类似。沿这条思路的证明见 [Woo65]。

\takeaway{$K$ 与 $KO$ 理论既可用分类空间 $BU\times\Z$、$BO\times\Z$ 的同伦类刻画，也可用向量丛的形式差来定义；两种定义等价，且分别具有 2-周期与 8-周期性。}

\section{$d+1$ 维的有能隙哈密顿量}

\subsection{凝聚态分析}

现在考虑 $d+1$ 维中的同一问题。也就是说，设有一个自由费米子系统，由形式如下的单粒子哈密顿量描述：
\begin{equation}
H=p_1\alpha_1+p_2\alpha_2+\cdots+p_d\alpha_d+h_0, \tag{11.48}
\end{equation}
其中与 $p$ 无关的部分 $h_0$ 具有复/实对称型 $i$。为使能谱具有形式
\begin{equation}
E^2=(p_1)^2+\cdots+(p_d)^2+m^2, \tag{11.49}
\end{equation}
我们希望各 $\alpha_k$ 彼此反对易，并且与 $h_0$ 反对易，于是多出额外的 Clifford 生成元，实际上把对称型移动了 $d$ 步。一个微妙之处是：$p_k$ 在时间反演下为奇，因此也希望 $\alpha_k$ 为奇，即 $T\alpha_k=-\alpha_kT$。这与我们上面考虑的费米子的行为相反，所以最终是在表 7 与表 8 中“向上”移位。这意味着 $d+1$ 维中复/实对称型 $i$ 的费米子自由有能隙哈密顿量分别由
\begin{equation}
K^{d-i}(\pt),\qquad KO^{d-i}(\pt) \tag{11.50}
\end{equation}
分类。
% [原文 p.169]
这给出了拓扑绝缘体与拓扑超导体的著名“周期表”：对拓扑绝缘体，电磁 $U(1)$ 未被破坏，因此只考虑不含反酉算子的对称类。对拓扑超导体，Cooper 对一般把电磁 $U(1)$ 破坏为 $\Z_2$，这可以通过反酉对称性算子的存在来实现 [Kit09, SRFL08, RSFL09]。见表 9。

\begin{center}
\footnotesize
\setlength{\tabcolsep}{3.4pt}
\begin{tabular}[t]{c||c|cc}
\toprule
name & $i\backslash d$ & 0 & 1\\
\midrule
A & 0 & $\Z$ & 0\\
AIII & 1 & 0 & $\Z$\\
\bottomrule
\end{tabular}
\hspace{10pt}
\begin{tabular}[t]{c||c|cccccccc|c}
\toprule
name & $i\backslash d$ & 0 & 1 & 2 & 3 & 4 & 5 & 6 & 7 & rel.\\
\midrule
AI   & 0 & $\Z$  & 0    & 0    & 0    & $\Z$ & 0    & $\Z_2$ & $\Z_2$ & \\
BDI  & 1 & $\Z_2$& $\Z$ & 0    & 0    & 0    & $\Z$ & 0      & $\Z_2$ & pin$^-$\\
D    & 2 & $\Z_2$& $\Z_2$& $\Z$ & 0    & 0    & 0    & $\Z$   & 0      & spin\\
DIII & 3 & 0    & $\Z_2$& $\Z_2$& $\Z$ & 0    & 0    & 0      & $\Z$   & pin$^+$\\
AII  & 4 & $\Z$  & 0    & $\Z_2$& $\Z_2$& $\Z$ & 0    & 0      & 0      & \\
CII  & 5 & 0    & $\Z$ & 0    & $\Z_2$& $\Z_2$& $\Z$ & 0      & 0      & \\
CI   & 6 & 0    & 0    & $\Z$ & 0    & $\Z_2$& $\Z_2$& $\Z$  & 0      & \\
C    & 7 & 0    & 0    & 0    & $\Z$ & 0    & $\Z_2$& $\Z_2$ & $\Z$  & \\
\bottomrule
\end{tabular}

\smallskip
{\small 表 9：拓扑绝缘体与拓扑超导体的周期表。}
\end{center}

这些构造统一了以往为拓扑绝缘体与超导体设计出的各种不变量。例如，在 $2+1$ 维中，整数量子霍尔效应的 Chern 数或等价的霍尔电导对应于对称类 A 的 $\Z$；热霍尔电导对应于对称类 D 的 $\Z$；而类 DIII 的 $\Z_2$ 分类由 Kane--Mele 不变量完成。

\subsection{相对论性分析}

相对论性费米子也遵循同样的分类。考虑 $D=d+1$ 维中在 Lorentz 群单位元连通分支下不变相对论性费米子。我们在关于旋量的一节中已经看到，它们是 Clifford 代数 $\mathrm{Cliff}_{\R}(d,0)$ 的不分次表示，并且
\begin{equation}
\{\alpha^i,\alpha^j\}=+2\delta^{ij} \tag{11.51}
\end{equation}
其中当我们使用 Majorana 基时 $\alpha^i=\gamma^0\gamma^i$。于是有质量 Dirac 方程具有形式
\begin{equation}
\partial_t\psi=\sum_i\alpha^i\partial_i\psi+m\beta\psi, \tag{11.52}
\end{equation}
其中 $m\beta$ 是实质量矩阵。为使它有 $E^2=\sum(p_i)^2+m^2$ 的色散关系，$\beta$ 需要与各 $\alpha^i$ 反对易，并满足 $\beta^2=-1$。因此我们考虑的是把 $\mathrm{Cliff}_{\R}(d,0)$ 的不分次表示延拓为 $\mathrm{Cliff}_{\R}(d,1)$ 的表示。等价地，是从 $\mathrm{Cliff}_{\R}(0,d-2)$ 到 $\mathrm{Cliff}_{\R}(1,d-2)$ 的延拓，或等价地，从 $\mathrm{Cliff}_{\R}(8n-(d-2),0)$ 到 $\mathrm{Cliff}_{\R}(8n-(d-2)+1,0)$ 的延拓。这由 $KO^{d-2}(\pt)$ 分类，所以这是对称型 D。

% [原文 p.170]
看出这是对称型 D 的另一种办法，是把式 (11.52) 与一般哈密顿量 (11.48) 比较。此时式 (11.48) 中的 $h_0$ 对应于式 (11.52) 中的 $i\beta$，即 (11.48) 中的 $h_0$ 被迫为纯虚，因此存在反酉对称性 $\tilde T$，满足 $\tilde Th_0=-h_0\tilde T$，而这正是对称类 D 的定义性质。这个反酉对称性 $\tilde T$ 之所以自动存在，是因为系统具有相对论性；它体现了相对论量子场论中 CPT 对称性总是存在这一事实。

我们也可以考虑带宇称对称性的相对论情形。如上所述，此时的 Dirac 方程为
\begin{equation}
\gamma_0\partial_t\psi-\gamma_i\partial_i\psi+m\gamma_m\psi=0, \tag{11.53}
\end{equation}
其中 $\gamma_0$ 与 $\gamma_m$ 类型相同，而各 $\gamma_i$ 类型相反。具体地说，在 pin$^+$ 情形考虑从 $\mathrm{Cliff}_{\R}(d,1)$ 到 $\mathrm{Cliff}_{\R}(d,2)$ 的延拓问题，而在 pin$^-$ 情形考虑从 $\mathrm{Cliff}_{\R}(1,d)$ 到 $\mathrm{Cliff}_{\R}(2,d)$ 的延拓问题。因此 pin$^+$ 情形由 $KO^{d-3}(\pt)$ 分类，属对称型 DIII；而 pin$^-$ 情形由 $KO^{d-1}(\pt)$ 分类，属对称型 BDI。

这意味着，例如 $2+1$ 维中大质量费米子的分类对 pin$^-$ 是平凡的，对 spin 是 $\Z$，对 pin$^+$ 是 $\Z_2$。这与我们在 §11.1 中所作的显式分析一致。

\takeaway{$d+1$ 维中，$\alpha_k$ 与 $h_0$ 反对易相当于多出 $d$ 个 Clifford 生成元，分类由 $K^{d-i}(\pt)$、$KO^{d-i}(\pt)$ 给出，即表 9 的周期表；相对论情形用 $\mathrm{Cliff}_{\R}$ 的延拓问题得到同样结果。}

\chapterreview{§11.1 用 $1+1$ 维边界与 $2+1$ 维体的两个例子说明：无额外不变性时边界无能隙自由度按 $\Z$ 计量，带 $T^2=(-1)^F$ 时按 $\Z_2$ 计量。§11.2 由随机矩阵论引出十重对称类，包括三个经典类、手征类与 BdG 类，并以 Cartan 的十族对称空间给出其酉算子空间（表 7、表 8）。§11.3 说明每加一个费米子算子，对称类就移位一步：复类 2-周期，实类 8-周期，组合系统的类按 $i+j$ 相加。§11.4 把 $U=e^{iH}$ 中的 $H$ 形变到本征值 $\pm1$，用直接分析与 Clifford 代数两种方法把有能隙哈密顿量的分类化为分类空间问题。§11.5 给出 $K$ 与 $KO$ 理论的两种定义（同伦类与向量丛形式差），并说明其 2-周期与 8-周期性。§11.6 把上述结果推广到 $d+1$ 维，得到拓扑绝缘体与拓扑超导体的周期表（表 9）。}
